% ============================================================================
% SPANISCH LANGENTWURF - DS 5: ASSESSMENT (SONSTIGE NOTE)
% ============================================================================
% Jahrgangsübergreifend Spätbeginner (Kl. 10, 11, 12)
% Datum: Mo 19.01.2026
% ============================================================================

\documentclass[11pt,a4paper]{article}

% ============================================================================
% PACKAGES
% ============================================================================
\usepackage[utf8]{inputenc}
\usepackage[T1]{fontenc}
\usepackage[ngerman]{babel}
\usepackage{geometry}
\usepackage{fancyhdr}
\usepackage{lastpage}
\usepackage{xcolor}
\usepackage{tcolorbox}
\usepackage{enumitem}
\usepackage{tabularx}
\usepackage{longtable}
\usepackage{array}
\usepackage{booktabs}
\usepackage{hyperref}
\usepackage{ifthen}
\usepackage{amssymb}
\usepackage{graphicx}
\usepackage{multirow}

% ============================================================================
% KONFIGURATION
% ============================================================================
\newcommand{\plantier}{1}
\newcommand{\fachid}{spanisch}
\newcommand{\fach}{Spanisch}
\newcommand{\fachfarbe}{c41e3a}

% --- Metadaten ---
\newcommand{\jahrgangsstufe}{10--12 (Spätbeginner)}
\newcommand{\schulart}{Gesamtschule -- Jahrgangsübergreifend}
\newcommand{\stundenthema}{Assessment -- Sonstige Note}
\newcommand{\reihenthema}{Mi barrio / Mi profesión ideal}
\newcommand{\stundennummer}{5}
\newcommand{\reihenstunden}{8}
\newcommand{\unterrichtsdatum}{19.01.2026}
\newcommand{\unterrichtszeit}{--:-- -- --:--}
\newcommand{\raum}{---}

% --- Lehrkraft ---
\newcommand{\lehrkraft}{N.N.}
\newcommand{\ausbildungsstand}{Lehrkraft}

% --- Lerngruppe ---
\newcommand{\lerngruppengroesse}{10}
\newcommand{\maennlich}{--}
\newcommand{\weiblich}{--}
\newcommand{\divers}{--}

% --- Kompetenzen/Niveau ---
\newcommand{\gerniveau}{A1--A2}
\newcommand{\lehrwerk}{A\_tope.com}
\newcommand{\lehrwerkseite}{Lección 3--4}

% ============================================================================
% LAYOUT
% ============================================================================
\geometry{left=2.5cm, right=2.5cm, top=2.5cm, bottom=2.5cm}
\definecolor{fach}{HTML}{\fachfarbe}
\definecolor{gray}{HTML}{666666}
\definecolor{lightgray}{HTML}{f5f5f5}
\definecolor{successgreen}{HTML}{2a7a4b}
\definecolor{warningorange}{HTML}{b35c00}

\tcbuselibrary{skins,breakable}

\newtcolorbox{infobox}[1][]{
    colback=lightgray,
    colframe=fach,
    fonttitle=\bfseries,
    title=#1,
    breakable
}

\newtcolorbox{scorebox}{
    colback=white,
    colframe=fach,
    fonttitle=\bfseries,
    title=Didaktik-Score,
    breakable
}

\newtcolorbox{warnbox}{
    colback=white,
    colframe=warningorange,
    fonttitle=\bfseries,
    title=Organisation,
    breakable
}

% Kopf-/Fußzeile
\pagestyle{fancy}
\fancyhf{}
\fancyhead[L]{\small\textcolor{gray}{\fach{} -- Jg. \jahrgangsstufe{} -- \stundenthema}}
\fancyhead[R]{\small\textcolor{gray}{Langentwurf Tier 1}}
\fancyfoot[C]{\small\thepage{} / \pageref{LastPage}}
\renewcommand{\headrulewidth}{0.4pt}
\renewcommand{\footrulewidth}{0pt}

% ============================================================================
% DOKUMENT
% ============================================================================
\begin{document}

% ============================================================================
% DECKBLATT
% ============================================================================
\begin{titlepage}
    \centering
    \vspace*{2cm}

    {\large\textcolor{gray}{\schulart}}\\[2cm]

    {\Huge\textcolor{fach}{\textbf{Langentwurf}}}\\[0.5cm]
    {\large Tier 1 -- Vollständig}\\[2cm]

    {\LARGE\textbf{\stundenthema}}\\[0.5cm]
    {\large Mündliche Prüfung zur Sonstigen Note}\\[0.3cm]
    {\Large\textit{Gruppe A: Mi barrio | Gruppe B: Mi profesión ideal}}\\[2cm]

    \begin{tabular}{rl}
        \textbf{Fach:} & \fach{} (\gerniveau) \\[0.3cm]
        \textbf{Klasse:} & \jahrgangsstufe{} \\[0.3cm]
        \textbf{Lehrwerk:} & \lehrwerk, \lehrwerkseite \\[0.3cm]
        \textbf{Datum:} & \underline{\hspace{1cm}Mo, 19.01.2026\hspace{1cm}} \\[0.3cm]
        \textbf{Zeit:} & \underline{\hspace{3cm}} (Doppelstunde) \\[0.3cm]
        \textbf{Raum:} & \underline{\hspace{3cm}} \\[0.3cm]
    \end{tabular}

    \vfill

    {\large Lehrkraft: \lehrkraft}\\[0.3cm]
    {\normalsize\textcolor{gray}{\ausbildungsstand}}\\[1cm]

\end{titlepage}

% ============================================================================
% INHALTSVERZEICHNIS
% ============================================================================
\tableofcontents
\newpage

% ============================================================================
% 1. BEDINGUNGSANALYSE
% ============================================================================
\section{Bedingungsanalyse}

\subsection{Lerngruppe}
\begin{tabular}{ll}
    Kursgröße: & 10 SuS (jahrgangsübergreifend) \\
    Jahrgangsstufen: & 10, 11, 12 \\
    Kurstyp: & Spätbeginner (2./3. Fremdsprache) \\
    Sprachniveau: & A1--A2 (GER) \\
\end{tabular}

\subsection{Prüfungsgruppen}

\begin{warnbox}
\textbf{Parallele Prüfungsorganisation:}
\begin{itemize}
    \item \textbf{Gruppe A (8 SuS):} Mündliche Partnerprüfung ``Mi barrio''
    \item \textbf{Gruppe B (2 SuS):} Kurzvortrag ``Mi profesión ideal''
\end{itemize}
Gruppe B präsentiert, während Gruppe A Partnergespräche führt (Rotation).
\end{warnbox}

\subsubsection{Gruppe A: Partnerprüfung ``Mi barrio'' (4 Paare)}

\begin{tabular}{|l|l|l|p{5cm}|}
\hline
\textbf{Paar} & \textbf{Partner 1} & \textbf{Partner 2} & \textbf{Niveaupaarung} \\
\hline
1 & E.H. & A.S. & stark + schwach \\
\hline
2 & L.N. & A.A. & stark + unsicher \\
\hline
3 & C.P. & V.S. & beide ruhig \\
\hline
4 & L.K. & H.P. & beide nachholen (flexibel) \\
\hline
\end{tabular}

\vspace{0.5em}
\textbf{Format:} 2 SuS führen Dialog, 3. beobachtet (optionale Rotation)\\
\textbf{Dauer:} 3--4 Min pro Paar + Vorbereitung

\subsubsection{Gruppe B: Kurzvortrag ``Mi profesión ideal'' (2 SuS)}

\begin{tabular}{|l|p{8cm}|}
\hline
\textbf{SuS} & \textbf{Hinweise} \\
\hline
[Name 1] & Einzelvortrag 3--4 Min \\
\hline
[Name 2] & Einzelvortrag 3--4 Min \\
\hline
\end{tabular}

\vspace{0.5em}
\textbf{Format:} Freies Sprechen mit Stichpunktkarten

\subsection{Besonderheiten}
\begin{itemize}
    \item \textbf{Jahrgangsübergreifend:} Unterschiedliche Vorerfahrungen (LJ1/LJ2/LJ3)
    \item \textbf{Prüfungsangst:} Bei ruhigen SuS (C.P., V.S.) durch Partnersituation reduzieren
    \item \textbf{Nachholer:} L.K., H.P. flexibel einplanen (ggf. eigenes Paar oder Dreiergruppe)
\end{itemize}

% ============================================================================
% 2. SACHANALYSE
% ============================================================================
\section{Sachanalyse}

\subsection{Gruppe A: ``Mi barrio'' -- Sprachliche Strukturen}

\begin{tabular}{|l|p{6cm}|p{5cm}|}
\hline
\textbf{Bereich} & \textbf{Struktur} & \textbf{Beispiele} \\
\hline
\textbf{Verben} & ser/estar/hay (Unterscheidung) & Mi barrio es tranquilo. / Hay una panadería. / La iglesia está en el centro. \\
\hline
\textbf{Fragen} & ¿Hay...? / ¿Cómo es...? / ¿Adónde vas...? & ¿Hay un parque cerca? / ¿Cómo es tu calle? \\
\hline
\textbf{Wortschatz} & Orte im Stadtviertel & el supermercado, la farmacia, el parque, la parada de autobús, la iglesia \\
\hline
\textbf{Adjektive} & Beschreibung & tranquilo, ruidoso, grande, pequeño, moderno, antiguo \\
\hline
\end{tabular}

\subsection{Gruppe B: ``Mi profesión ideal'' -- Sprachliche Strukturen}

\begin{tabular}{|l|p{6cm}|p{5cm}|}
\hline
\textbf{Bereich} & \textbf{Struktur} & \textbf{Beispiele} \\
\hline
\textbf{Berufswortschatz} & profesiones & médico/a, profesor/a, ingeniero/a, abogado/a, periodista \\
\hline
\textbf{Modalverben} & querer, poder, necesitar & Quiero ser... / Necesito estudiar... \\
\hline
\textbf{Konjunktionen} & porque, para, cuando & Quiero ser médico porque me gusta ayudar. \\
\hline
\textbf{Zukunft} & ir a + infinitivo & Voy a estudiar medicina. \\
\hline
\end{tabular}

\subsection{Kontrastive Betrachtung}
\begin{itemize}
    \item \textbf{ser/estar:} Keine direkte Entsprechung im Deutschen -- häufige Fehlerquelle
    \item \textbf{hay:} Unpersönliches ``es gibt'' -- einfacher als ser/estar
    \item \textbf{Berufsartikel:} Soy médico (ohne Artikel) vs. Dt. ``Ich bin Arzt''
\end{itemize}

% ============================================================================
% 3. DIDAKTISCHE ANALYSE
% ============================================================================
\section{Didaktische Analyse}

\subsection{Stellung in der Sequenz}
Dies ist die \textbf{5. Doppelstunde} der Sequenz (8 Doppelstunden gesamt).

\textbf{Assessment-Stunde:} Überprüfung der mündlichen Kompetenz für ``Sonstige Note''.

\subsection{Prüfungskompetenzen (Rahmenplan MV)}

\textbf{Gruppe A:}
\begin{itemize}
    \item \textbf{Sprechen -- an Gesprächen teilnehmen:} Einfache Dialoge über vertraute Themen führen
    \item \textbf{Wortschatz:} Orte und Beschreibungen im Stadtviertel
    \item \textbf{Grammatik:} Korrekte Verwendung von ser/estar/hay
\end{itemize}

\textbf{Gruppe B:}
\begin{itemize}
    \item \textbf{Sprechen -- zusammenhängendes Sprechen:} Kurze Präsentation zu persönlichem Thema
    \item \textbf{Argumentation:} Begründungen mit porque, para formulieren
    \item \textbf{Kohärenz:} Strukturierter Aufbau des Vortrags
\end{itemize}

\subsection{Legitimation}
\textbf{Rahmenplan Spanisch MV (Sek I/II):}
\begin{itemize}
    \item Mündliche Kommunikationsfähigkeit als Kernkompetenz
    \item Dialogisches und monologisches Sprechen gleichermaßen fördern
    \item Leistungsbewertung durch mündliche Prüfungsformate
\end{itemize}

% ============================================================================
% 4. STUNDENZIELE
% ============================================================================
\section{Stundenziele}

\subsection{Grobziel}
Die SuS demonstrieren ihre \textbf{mündliche Kommunikationskompetenz} (dialogisch und monologisch), indem sie in strukturierten Prüfungssituationen auf Spanisch kommunizieren.

\subsection{Feinziele (Gruppe A)}
\begin{tabular}{|c|p{7cm}|c|c|}
\hline
\textbf{Nr.} & \textbf{Die SuS können...} & \textbf{AFB} & \textbf{Punkte} \\
\hline
FZ1 & ihr Stadtviertel mit ser/estar/hay beschreiben & II & 5,4 \\
\hline
FZ2 & passende Fragen zum Stadtviertel formulieren & II & 4,5 \\
\hline
FZ3 & auf Fragen grammatisch korrekt antworten & II & 3,6 \\
\hline
FZ4 & flüssig und verständlich kommunizieren & II & 4,5 \\
\hline
\end{tabular}

\subsection{Feinziele (Gruppe B)}
\begin{tabular}{|c|p{7cm}|c|c|}
\hline
\textbf{Nr.} & \textbf{Die SuS können...} & \textbf{AFB} & \textbf{Punkte} \\
\hline
FZ1 & ihren Berufswunsch mit Begründung präsentieren & II & 5,4 \\
\hline
FZ2 & notwendige Fähigkeiten benennen & I & 4,5 \\
\hline
FZ3 & den Ausbildungsweg strukturiert darstellen & II & 3,6 \\
\hline
FZ4 & flüssig und mit guter Aussprache vortragen & II & 4,5 \\
\hline
\end{tabular}

% ============================================================================
% 5. METHODISCHE ANALYSE
% ============================================================================
\section{Methodische Analyse}

\subsection{Prüfungsformate}

\subsubsection{Gruppe A: Partnerprüfung (dialogisch)}
\begin{itemize}
    \item \textbf{Aufgabe:}
    \begin{enumerate}
        \item Partner A beschreibt sein Stadtviertel (1,5--2 Min)
        \item Partner B stellt 3 Fragen (¿Hay...? ¿Cómo es...? ¿Adónde vas...?)
        \item Rollentausch
    \end{enumerate}
    \item \textbf{Dauer:} 3--4 Min pro Paar
    \item \textbf{Beobachter:} Optional ein dritter SuS (nicht bewertet)
\end{itemize}

\subsubsection{Gruppe B: Kurzvortrag (monologisch)}
\begin{itemize}
    \item \textbf{Inhalt:}
    \begin{enumerate}
        \item Welchen Beruf möchtest du? Warum?
        \item Welche Fähigkeiten brauchst du?
        \item Wie ist der Weg dorthin?
    \end{enumerate}
    \item \textbf{Dauer:} 3--4 Min pro Person
    \item \textbf{Hilfsmittel:} Stichpunktkarten erlaubt (max. 5 Stichpunkte)
\end{itemize}

\subsection{Bewertung}

\textbf{Notenskala:} 18 Punkte gesamt $\rightarrow$ 14-NP-System (Sek I)

\begin{tabular}{|c|c|c|c|c|c|c|c|c|c|c|c|c|c|c|}
\hline
\textbf{NP} & 14 & 13 & 12 & 11 & 10 & 9 & 8 & 7 & 6 & 5 & 4 & 3 & 2 & 1 \\
\hline
\textbf{\%} & 100 & 96 & 91 & 86 & 80 & 73 & 67 & 60 & 53 & 47 & 40 & 33 & 27 & 20 \\
\hline
\textbf{Pkt.} & 18 & 17,3 & 16,4 & 15,5 & 14,4 & 13,1 & 12 & 10,8 & 9,5 & 8,5 & 7,2 & 5,9 & 4,9 & 3,6 \\
\hline
\end{tabular}

\subsection{Medien}
\begin{itemize}
    \item Bewertungsbogen Gruppe A (Anlage 1)
    \item Bewertungsbogen Gruppe B (Anlage 2)
    \item Aufgabenkarten für SuS
    \item Stoppuhr
    \item Sitzordnung für Prüfungsrotation
\end{itemize}

% ============================================================================
% 6. VERLAUFSPLANUNG
% ============================================================================
\section{Verlaufsplanung}

\begin{warnbox}
\textbf{Parallelbetrieb:}
\begin{itemize}
    \item Lehrkraft prüft Gruppe B (Vorträge) am Lehrertisch
    \item Gruppe A bereitet sich vor / führt Partnergespräche im hinteren Bereich
    \item Nach Gruppe B: Lehrkraft wechselt zu Gruppe A
\end{itemize}
\end{warnbox}

\vspace{1em}
\textbf{1. Stunde (45 Min.)}

\begin{longtable}{|p{1.2cm}|p{1.5cm}|p{4cm}|p{4cm}|p{1.5cm}|p{1.5cm}|}
\hline
\textbf{Zeit} & \textbf{Phase} & \textbf{Lehrerhandeln} & \textbf{Schülerhandeln} & \textbf{SF} & \textbf{Medien} \\
\hline
\endhead

0--5 & Orga & Begrüßung, Ablauf erklären, Gruppen zuweisen, Sitzordnung & Plätze einnehmen, Materialien bereit & Plenum & Tafel \\
\hline

5--10 & Vorber. & Aufgabenkarten ausgeben, Fragen klären & Aufgaben lesen, Stichpunkte notieren & EA & Aufg.-karten \\
\hline

\multicolumn{6}{|c|}{\textbf{Gruppe B: Vorträge (Lehrkraft am Lehrertisch)}} \\
\hline

10--15 & B1 & [SuS 1] prüfen, Bewertung notieren & [SuS 1] präsentiert ``Mi profesión ideal'' & Einzel & Bew.-bogen B \\
\hline

15--20 & B2 & [SuS 2] prüfen, Bewertung notieren & [SuS 2] präsentiert ``Mi profesión ideal'' & Einzel & Bew.-bogen B \\
\hline

\multicolumn{6}{|c|}{\textbf{Gruppe A: Vorbereitung (parallel, selbstständig)}} \\
\hline

10--20 & Vorber. & (Lehrkraft bei Gruppe B) & Paare bereiten Dialog vor, Stichpunkte & PA & Aufg.-karten \\
\hline

\multicolumn{6}{|c|}{\textbf{Gruppe A: Partnerprüfungen (Lehrkraft wechselt)}} \\
\hline

20--27 & A1 & Paar 1 (E.H. + A.S.) prüfen & Partnerdialog ``Mi barrio'' & PA & Bew.-bogen A \\
\hline

27--34 & A2 & Paar 2 (L.N. + A.A.) prüfen & Partnerdialog ``Mi barrio'' & PA & Bew.-bogen A \\
\hline

34--41 & A3 & Paar 3 (C.P. + V.S.) prüfen & Partnerdialog ``Mi barrio'' & PA & Bew.-bogen A \\
\hline

41--45 & Puffer & Puffer / Organisation & Wartende: Reflexionsbogen ausfüllen & EA & Refl.-bogen \\
\hline
\end{longtable}

\vspace{1em}

\textbf{2. Stunde (45 Min.)}

\begin{longtable}{|p{1.2cm}|p{1.5cm}|p{4cm}|p{4cm}|p{1.5cm}|p{1.5cm}|}
\hline
\textbf{Zeit} & \textbf{Phase} & \textbf{Lehrerhandeln} & \textbf{Schülerhandeln} & \textbf{SF} & \textbf{Medien} \\
\hline
\endhead

0--7 & A4 & Paar 4 (L.K. + H.P.) prüfen & Partnerdialog ``Mi barrio'' & PA & Bew.-bogen A \\
\hline

7--15 & Puffer & Eventuelle Nachholprüfungen / Wiederholungen & Wartende: Selbstreflexion & EA & -- \\
\hline

\multicolumn{6}{|c|}{\textbf{Übergang: Lehrkraft bei allen SuS}} \\
\hline

15--25 & Reflexion & Kurze Rückmeldung zu Stärken (keine Noten!) & Zuhören, Fragen stellen & Plenum & -- \\
\hline

25--35 & Übung & Vertiefungsübung: ser/estar/hay Wiederholung & Arbeitsblatt bearbeiten & EA & AB \\
\hline

35--42 & Sicherung & Lösungen besprechen, Fehlertypen erläutern & Vergleichen, korrigieren & Plenum & Tafel \\
\hline

42--45 & Abschluss & Ausblick nächste Stunde, Notenbekanntgabe & Fragen, Aufräumen & Plenum & -- \\
\hline
\end{longtable}

\textbf{Legende:} EA = Einzelarbeit, PA = Partnerarbeit, SF = Sozialform

% ============================================================================
% 7. ERWARTETE SCHWIERIGKEITEN
% ============================================================================
\section{Erwartete Schwierigkeiten}

\begin{tabular}{|p{4.5cm}|p{8cm}|}
\hline
\textbf{Schwierigkeit} & \textbf{Reaktion} \\
\hline
\textbf{Prüfungsangst} (bes. C.P., V.S.) & Partnersituation nutzen, ermutigen, Zeit geben \\
\hline
\textbf{ser/estar/hay-Verwechslung} & In Bewertung differenzieren: Versuche würdigen, nicht nur Fehler zählen \\
\hline
\textbf{Zeitüberschreitung} & Stoppuhr sichtbar, sanfter Hinweis bei 3:30 Min \\
\hline
\textbf{Partner spricht nicht} & Lehrkraft springt als Gesprächspartner ein \\
\hline
\textbf{Nachholer fehlen} & Dreiergruppe bilden oder Einzelprüfung (Lehrkraft als Partner) \\
\hline
\textbf{Unruhe bei Wartenden} & Reflexionsbogen ausfüllen, leise Vorbereitung \\
\hline
\end{tabular}

% ============================================================================
% 8. HAUSAUFGABE
% ============================================================================
\section{Hausaufgabe}

\begin{itemize}
    \item \textbf{Freiwillig:} Reflexion zur eigenen Prüfung (3 Sätze auf Spanisch)
    \begin{itemize}
        \item ¿Qué he hecho bien?
        \item ¿Qué puedo mejorar?
        \item ¿Qué quiero practicar más?
    \end{itemize}
    \item \textbf{Vorbereitung DS 6:} Vokabeln Thema [nächstes Thema] wiederholen
\end{itemize}

% ============================================================================
% 9. ANLAGEN
% ============================================================================
\section{Anlagen}

\begin{enumerate}
    \item Bewertungsbogen Gruppe A: Partnerprüfung ``Mi barrio''
    \item Bewertungsbogen Gruppe B: Kurzvortrag ``Mi profesión ideal''
    \item Aufgabenkarten für SuS (Gruppe A und B)
    \item Selbstreflexionsbogen
    \item Notenschlüssel 18 Punkte $\rightarrow$ 14-NP-System
\end{enumerate}

\newpage

% ============================================================================
% ANLAGE 1: BEWERTUNGSBOGEN GRUPPE A
% ============================================================================
\section*{Anlage 1: Bewertungsbogen Gruppe A -- Partnerprüfung ``Mi barrio''}
\addcontentsline{toc}{section}{Anlage 1: Bewertungsbogen Gruppe A}

\begin{center}
\textbf{\large Bewertungsbogen Mündliche Prüfung -- Mi barrio}\\[0.5em]
Datum: \underline{\hspace{2cm}} \hspace{2cm} Paar Nr.: \underline{\hspace{1cm}}
\end{center}

\vspace{0.5em}

\begin{tabular}{|l|c|c|}
\hline
\textbf{Partner A:} & \underline{\hspace{5cm}} & \textbf{Klasse:} \underline{\hspace{2cm}} \\
\hline
\textbf{Partner B:} & \underline{\hspace{5cm}} & \textbf{Klasse:} \underline{\hspace{2cm}} \\
\hline
\end{tabular}

\vspace{1em}

\textbf{Aufgabe:}
\begin{enumerate}
    \item Partner A beschreibt sein Stadtviertel (1,5--2 Min)
    \item Partner B stellt 3 Fragen: ¿Hay...? / ¿Cómo es...? / ¿Adónde vas...?
    \item Rollentausch
\end{enumerate}

\vspace{1em}

\subsection*{Bewertungskriterien}

\begin{tabular}{|p{4cm}|c|p{5cm}|c|c|}
\hline
\textbf{Kriterium} & \textbf{Gew.} & \textbf{Deskriptor} & \textbf{A} & \textbf{B} \\
\hline
\hline
\multirow{4}{*}{\textbf{1. ser/estar/hay}} & \multirow{4}{*}{30\%} & 5--6: Korrekte Verwendung, kaum Fehler & & \\
& & 3--4: Überwiegend korrekt, einige Fehler & & \\
& & 1--2: Häufige Verwechslung & & \\
& & 0: Keine Unterscheidung erkennbar & & \\
\cline{3-5}
& & \textbf{Punkte (max. 5,4):} & /5,4 & /5,4 \\
\hline
\hline
\multirow{4}{*}{\textbf{2. Wortschatz}} & \multirow{4}{*}{25\%} & 4--5: Reichhaltig, passende Orte/Adjektive & & \\
\textbf{(Stadtviertel/Orte)} & & 2--3: Grundwortschatz vorhanden & & \\
& & 1: Sehr eingeschränkt & & \\
& & 0: Unzureichend & & \\
\cline{3-5}
& & \textbf{Punkte (max. 4,5):} & /4,5 & /4,5 \\
\hline
\hline
\multirow{4}{*}{\textbf{3. Grammatik}} & \multirow{4}{*}{20\%} & 3--4: Weitgehend korrekte Strukturen & & \\
\textbf{(Korrektheit)} & & 2: Grundstrukturen vorhanden, Fehler & & \\
& & 1: Viele Fehler, aber verständlich & & \\
& & 0: Fehler behindern Verständnis & & \\
\cline{3-5}
& & \textbf{Punkte (max. 3,6):} & /3,6 & /3,6 \\
\hline
\hline
\multirow{4}{*}{\textbf{4. Flüssigkeit}} & \multirow{4}{*}{25\%} & 4--5: Flüssig, gute Aussprache, interaktiv & & \\
\textbf{(+ Verständlichkeit)} & & 2--3: Stockend, aber verständlich & & \\
& & 1: Häufige Pausen, schwer verständlich & & \\
& & 0: Kommunikation bricht ab & & \\
\cline{3-5}
& & \textbf{Punkte (max. 4,5):} & /4,5 & /4,5 \\
\hline
\hline
\multicolumn{3}{|r|}{\textbf{SUMME (max. 18):}} & \textbf{/18} & \textbf{/18} \\
\hline
\multicolumn{3}{|r|}{\textbf{NOTE (NP):}} & & \\
\hline
\end{tabular}

\vspace{1em}

\textbf{Notizen / Beobachtungen:}

\vspace{4cm}

\hrule

\vspace{0.5em}
Unterschrift Lehrkraft: \underline{\hspace{6cm}}

\newpage

% ============================================================================
% ANLAGE 2: BEWERTUNGSBOGEN GRUPPE B
% ============================================================================
\section*{Anlage 2: Bewertungsbogen Gruppe B -- Kurzvortrag ``Mi profesión ideal''}
\addcontentsline{toc}{section}{Anlage 2: Bewertungsbogen Gruppe B}

\begin{center}
\textbf{\large Bewertungsbogen Mündliche Prüfung -- Mi profesión ideal}\\[0.5em]
Datum: \underline{\hspace{2cm}}
\end{center}

\vspace{0.5em}

\begin{tabular}{|l|l|}
\hline
\textbf{Name:} & \underline{\hspace{7cm}} \\
\hline
\textbf{Klasse:} & \underline{\hspace{7cm}} \\
\hline
\textbf{Dauer:} & \underline{\hspace{3cm}} Min (Ziel: 3--4 Min) \\
\hline
\end{tabular}

\vspace{1em}

\textbf{Aufgabe:} Kurzvortrag (3--4 Min) zu:
\begin{enumerate}
    \item Welchen Beruf möchtest du? Warum?
    \item Welche Fähigkeiten brauchst du?
    \item Wie ist der Weg dorthin?
\end{enumerate}

\vspace{1em}

\subsection*{Bewertungskriterien}

\begin{tabular}{|p{4cm}|c|p{6cm}|c|}
\hline
\textbf{Kriterium} & \textbf{Gew.} & \textbf{Deskriptor} & \textbf{Punkte} \\
\hline
\hline
\multirow{4}{*}{\textbf{1. Argumentation}} & \multirow{4}{*}{30\%} & 5--6: Klare Begründungen, porque/para korrekt & \\
\textbf{(+ Konjunktionen)} & & 3--4: Begründungen vorhanden, teils fehlerhaft & \\
& & 1--2: Wenig Begründungen, einfache Sätze & \\
& & 0: Keine erkennbare Argumentation & \\
\cline{3-4}
& & \textbf{Punkte (max. 5,4):} & /5,4 \\
\hline
\hline
\multirow{4}{*}{\textbf{2. Fachvokabular}} & \multirow{4}{*}{25\%} & 4--5: Vielfältig (Berufe, Bildung, Fähigkeiten) & \\
\textbf{(Berufe/Bildung)} & & 2--3: Grundwortschatz vorhanden & \\
& & 1: Sehr eingeschränkt & \\
& & 0: Unzureichend für Thema & \\
\cline{3-4}
& & \textbf{Punkte (max. 4,5):} & /4,5 \\
\hline
\hline
\multirow{4}{*}{\textbf{3. Struktur}} & \multirow{4}{*}{20\%} & 3--4: Klarer Aufbau (Einleitung, Hauptteil, Schluss) & \\
\textbf{(+ Kohärenz)} & & 2: Erkennbare Gliederung, teils sprunghaft & \\
& & 1: Wenig Struktur, aber verständlich & \\
& & 0: Zusammenhanglose Einzelsätze & \\
\cline{3-4}
& & \textbf{Punkte (max. 3,6):} & /3,6 \\
\hline
\hline
\multirow{4}{*}{\textbf{4. Flüssigkeit}} & \multirow{4}{*}{25\%} & 4--5: Flüssig, gute Aussprache, frei vorgetragen & \\
\textbf{(+ Aussprache)} & & 2--3: Stockend, aber verständlich, Zettel genutzt & \\
& & 1: Häufige Pausen, stark abgelesen & \\
& & 0: Kommunikation kaum möglich & \\
\cline{3-4}
& & \textbf{Punkte (max. 4,5):} & /4,5 \\
\hline
\hline
\multicolumn{3}{|r|}{\textbf{SUMME (max. 18):}} & \textbf{/18} \\
\hline
\multicolumn{3}{|r|}{\textbf{NOTE (NP):}} & \\
\hline
\end{tabular}

\vspace{1em}

\textbf{Inhaltliche Stichpunkte (Vollständigkeit):}

\begin{tabular}{|l|c|}
\hline
Berufswunsch genannt & $\Box$ \\
\hline
Begründung (Warum?) & $\Box$ \\
\hline
Fähigkeiten genannt & $\Box$ \\
\hline
Ausbildungsweg beschrieben & $\Box$ \\
\hline
\end{tabular}

\vspace{1em}

\textbf{Notizen / Beobachtungen:}

\vspace{4cm}

\hrule

\vspace{0.5em}
Unterschrift Lehrkraft: \underline{\hspace{6cm}}

\newpage

% ============================================================================
% ANLAGE 3: AUFGABENKARTEN
% ============================================================================
\section*{Anlage 3: Aufgabenkarten für SuS}
\addcontentsline{toc}{section}{Anlage 3: Aufgabenkarten}

\vspace{1em}
\hrule
\vspace{1em}

\begin{center}
\fbox{\begin{minipage}{0.9\textwidth}
\begin{center}
\textbf{\large GRUPO A: MI BARRIO}\\
\textbf{Conversación con un compañero}
\end{center}

\vspace{1em}

\textbf{Tiempo:} 3--4 minutos (ambos juntos)

\vspace{0.5em}

\textbf{Tarea:}
\begin{enumerate}
    \item \textbf{Partner A:} Describe tu barrio (1,5--2 min)
    \begin{itemize}
        \item ¿Cómo es tu barrio? (tranquilo, ruidoso, grande...)
        \item ¿Qué hay en tu barrio? (tiendas, parques, iglesia...)
        \item ¿Dónde está tu casa?
    \end{itemize}
    \item \textbf{Partner B:} Haz 3 preguntas:
    \begin{itemize}
        \item ¿Hay...? (un supermercado, una farmacia...)
        \item ¿Cómo es...? (la calle, el parque...)
        \item ¿Adónde vas...? (para comprar, para jugar...)
    \end{itemize}
    \item \textbf{Cambio de roles}
\end{enumerate}

\vspace{0.5em}

\textbf{Vocabulario útil:}\\
hay -- es gibt | ser -- sein (Eigenschaft) | estar -- sein (Ort)\\
el supermercado, la farmacia, la panadería, el parque, la iglesia, la parada de autobús\\
tranquilo, ruidoso, grande, pequeño, moderno, antiguo, cerca, lejos
\end{minipage}}
\end{center}

\vspace{2em}
\hrule
\vspace{2em}

\begin{center}
\fbox{\begin{minipage}{0.9\textwidth}
\begin{center}
\textbf{\large GRUPO B: MI PROFESIÓN IDEAL}\\
\textbf{Presentación individual}
\end{center}

\vspace{1em}

\textbf{Tiempo:} 3--4 minutos

\vspace{0.5em}

\textbf{Tarea:} Presenta tu profesión ideal. Habla sobre:
\begin{enumerate}
    \item \textbf{Tu sueño:} ¿Qué quieres ser? ¿Por qué?
    \begin{itemize}
        \item Quiero ser... porque...
    \end{itemize}
    \item \textbf{Las habilidades:} ¿Qué necesitas saber/poder hacer?
    \begin{itemize}
        \item Necesito saber... / Tengo que poder...
    \end{itemize}
    \item \textbf{El camino:} ¿Cómo llegas a esta profesión?
    \begin{itemize}
        \item Primero voy a... / Después tengo que...
    \end{itemize}
\end{enumerate}

\vspace{0.5em}

\textbf{Vocabulario útil:}\\
médico/a, profesor/a, ingeniero/a, abogado/a, periodista, programador/a\\
estudiar, trabajar, aprender, practicar\\
la universidad, el instituto, las prácticas, el examen\\
porque, para, primero, después, finalmente
\end{minipage}}
\end{center}

\newpage

% ============================================================================
% ANLAGE 4: SELBSTREFLEXIONSBOGEN
% ============================================================================
\section*{Anlage 4: Selbstreflexionsbogen}
\addcontentsline{toc}{section}{Anlage 4: Selbstreflexionsbogen}

\begin{center}
\textbf{\large Reflexión -- Mi presentación oral}\\[0.5em]
Nombre: \underline{\hspace{5cm}} \hspace{1cm} Fecha: \underline{\hspace{3cm}}
\end{center}

\vspace{1em}

\textbf{1. Antes de la prueba:}

\begin{tabular}{|p{10cm}|c|c|c|}
\hline
& \textbf{Sí} & \textbf{Un poco} & \textbf{No} \\
\hline
Ich habe mich gut vorbereitet. & $\Box$ & $\Box$ & $\Box$ \\
\hline
Ich habe die Vokabeln gelernt. & $\Box$ & $\Box$ & $\Box$ \\
\hline
Ich habe laut geübt. & $\Box$ & $\Box$ & $\Box$ \\
\hline
\end{tabular}

\vspace{1em}

\textbf{2. Durante la prueba:}

\begin{tabular}{|p{10cm}|c|c|c|}
\hline
& \textbf{Sí} & \textbf{Un poco} & \textbf{No} \\
\hline
Ich konnte flüssig sprechen. & $\Box$ & $\Box$ & $\Box$ \\
\hline
Ich habe ser/estar/hay richtig benutzt. & $\Box$ & $\Box$ & $\Box$ \\
\hline
Ich habe meinen Partner verstanden. (Grupo A) & $\Box$ & $\Box$ & $\Box$ \\
\hline
Meine Struktur war klar. (Grupo B) & $\Box$ & $\Box$ & $\Box$ \\
\hline
\end{tabular}

\vspace{1em}

\textbf{3. ¿Qué he hecho bien?} (Was ist mir gut gelungen?)

\vspace{2cm}

\hrule

\vspace{1em}

\textbf{4. ¿Qué puedo mejorar?} (Was kann ich verbessern?)

\vspace{2cm}

\hrule

\vspace{1em}

\textbf{5. ¿Qué quiero practicar más?} (Was möchte ich mehr üben?)

\vspace{2cm}

\hrule

\newpage

% ============================================================================
% ANLAGE 5: NOTENSCHLÜSSEL
% ============================================================================
\section*{Anlage 5: Notenschlüssel 18 Punkte $\rightarrow$ 14-NP-System}
\addcontentsline{toc}{section}{Anlage 5: Notenschlüssel}

\begin{center}
\textbf{\large Umrechnungstabelle: Rohpunkte zu Notenpunkten}
\end{center}

\vspace{1em}

\begin{tabular}{|c|c|c|c|}
\hline
\textbf{Notenpunkte (NP)} & \textbf{Prozent} & \textbf{Rohpunkte (von 18)} & \textbf{Note (6-Skala)} \\
\hline
\hline
14 & 100\% & 18,0 & 1+ \\
\hline
13 & 96\% & 17,3 -- 17,9 & 1 \\
\hline
12 & 91\% & 16,4 -- 17,2 & 1-- \\
\hline
\hline
11 & 86\% & 15,5 -- 16,3 & 2+ \\
\hline
10 & 80\% & 14,4 -- 15,4 & 2 \\
\hline
9 & 73\% & 13,1 -- 14,3 & 2-- \\
\hline
\hline
8 & 67\% & 12,0 -- 13,0 & 3+ \\
\hline
7 & 60\% & 10,8 -- 11,9 & 3 \\
\hline
6 & 53\% & 9,5 -- 10,7 & 3-- \\
\hline
\hline
5 & 47\% & 8,5 -- 9,4 & 4+ \\
\hline
4 & 40\% & 7,2 -- 8,4 & 4 \\
\hline
3 & 33\% & 5,9 -- 7,1 & 4-- \\
\hline
\hline
2 & 27\% & 4,9 -- 5,8 & 5+ \\
\hline
1 & 20\% & 3,6 -- 4,8 & 5 \\
\hline
0 & $<$20\% & 0 -- 3,5 & 6 \\
\hline
\end{tabular}

\vspace{2em}

\textbf{Schnellübersicht:}

\begin{tabular}{|l|l|}
\hline
\textbf{Punkte} & \textbf{Tendenz} \\
\hline
16--18 & Sehr gut (1) \\
\hline
13--15,9 & Gut (2) \\
\hline
10--12,9 & Befriedigend (3) \\
\hline
7--9,9 & Ausreichend (4) \\
\hline
4--6,9 & Mangelhaft (5) \\
\hline
0--3,9 & Ungenügend (6) \\
\hline
\end{tabular}

% ============================================================================
% 10. DIDAKTIK-SCORE
% ============================================================================
\newpage
\section{Didaktik-Score}

\begin{scorebox}
\textbf{Bewertung der didaktischen Qualität nach 10 Dimensionen}\\
\textbf{Ziel-Score: $\geq$ 9.0 (Premium-Qualität)}

\vspace{1em}
\begin{tabular}{|c|l|c|c|p{4.5cm}|}
\hline
\textbf{\#} & \textbf{Dimension} & \textbf{Gew.} & \textbf{Score} & \textbf{Notiz} \\
\hline
1 & Differenzierung & 15\% & 9/10 & Niveau-Paarung, 2 Formate \\
\hline
2 & Taxonomie (AFB) & 10\% & 9/10 & AFB I--II (Prüfung) \\
\hline
3 & Lernziel-Alignment & 10\% & 10/10 & Kommunikativ, rahmenplankonform \\
\hline
4 & Aktivierung & 10\% & 10/10 & 100\% mündliche Produktion \\
\hline
5 & Scaffolding & 10\% & 9/10 & Aufgabenkarten, Stichpunkte \\
\hline
6 & Feedback-Qualität & 10\% & 9/10 & Bewertungsbogen, Reflexion \\
\hline
7 & Sprachliche Klarheit & 10\% & 9/10 & Klare Aufgabenstellung \\
\hline
8 & Motivation \& Relevanz & 10\% & 9/10 & Alltagsthemen, Berufsorientierung \\
\hline
9 & Inklusion (UDL) & 10\% & 8/10 & Partner unterstützen sich \\
\hline
10 & Praktikabilität & 5\% & 9/10 & Parallelbetrieb organisiert \\
\hline
\multicolumn{3}{|r|}{\textbf{GESAMT:}} & \textbf{9.1/10} & \\
\hline
\end{tabular}

\vspace{1em}
\textbf{Bewertungsschwellen:}
\begin{itemize}
    \item[\checkmark] \textcolor{successgreen}{$\geq$ 9.0: \textbf{FREIGABE} (Premium-Qualität)}
    \item[$\Box$] \textcolor{warningorange}{8.0--8.9: Iteration empfohlen}
    \item[$\Box$] 6.0--7.9: Überarbeitung erforderlich
    \item[$\Box$] $<$ 6.0: Grundlegende Neukonzeption
\end{itemize}
\end{scorebox}

\end{document}
